% ============================================================
% 分片：s11-appendix  对应原文 PDF 第 124--127 页（印刷页 124--127）
% 负责章节：§11 Appendix. L-packets for SL2
% 译者：trans-s07   审校状态：初稿
% ============================================================

% ==== 原文 PDF 第 124 页（印刷页 124）====
% [跨页接续: 原文第 124 页开头仍属 §10（由 s10-gerbes 分片负责），本分片自该页 §11 标题起]

\section{附录：\texorpdfstring{$SL_2$}{SL2} 的 L-包}

\begin{keytakeaway}
\textbf{核心物理图像：}本附录用 Whittaker 函数给出定理 7.1 的显式证明：$SL_2$ 的 L-包判据归结为 Whittaker 函数在一族特殊元素上是否恒为零，而后者由交换类域论给出的恒等式 (11.6) 控制。
\textbf{数学工具：}Jacquet--Langlands 的 Whittaker 函数 Fourier 变换、球 Hecke 代数、局部域上的加法特征标与留数、二次特征标 $\alpha: W_F\to\mathbb{Z}_2$、除子类的线性等价、$\langle M\rangle$ 的奇偶性与 Weil 的 $\langle\delta\rangle$ 恒偶定理。
\textbf{主要结论：}$\langle\delta+D\rangle$ 为奇数时 $f_{\tilde\sigma}$ 在 $\Bun^{\mathcal{O}(D)}_{SL_2}$ 的像上的限制为零，为偶数时该限制非零且相应的 Hecke 本征函数空间一维；由于 $\langle\delta\rangle$ 恒为偶数，定理 7.1 的 $\#S$ 奇偶判据得证。
\end{keytakeaway}

在本附录中，我们利用 Whittaker 函数概述定理 7.1 的一个证明。用 Whittaker 函数的 Fourier 变换构造 $GL_2$ 的自守函数，是由 H. Jacquet 与 R. Langlands [JL] 利用 A. Weil [W2] 的一个结果引入的。下文我们采用 [Dr2] 的表述与记号。

我们将固定 $C$ 上一个非零有理微分 $\omega$，并用 $\delta=\sum_x\delta_x[x]$ 表示它的零点与极点的除子。设 $\psi:\mathbb{F}_q\to\overline{\mathbb{Q}}_\ell^\times$ 是一个非平凡的加法特征标。它给出一个特征标 $\Psi:\mathbb{A}_F\to\overline{\mathbb{Q}}_\ell^\times$，由公式
\[
\Psi((f_x)) = \prod_{x\in C}\psi\left(\Tr_{k_x/\mathbb{F}_q}(\operatorname{Res}_x(f_x\omega))\right)
\]
定义。由留数公式，它在 $F\subset\mathbb{A}_F$ 上的值等于 1。因此 $\Psi$ 给出 $\mathbb{A}_F/F$ 的一个特征标。

设 $B\subset GL_2$ 是上三角矩阵构成的 Borel 子群。用 $V$ 表示局部常值函数 $f:GL_2(\mathbb{A}_F)\to\overline{\mathbb{Q}}_\ell$ 构成的空间，使得
\begin{itemize}[leftmargin=2.2em,itemsep=2pt,topsep=2pt]
\item[(1)] $f(\alpha x)=f(x)$ 对所有 $x\in GL_2(\mathbb{A}_F)$、$\alpha\in B(F)$ 成立；
\item[(2)] $\displaystyle\int_{\mathbb{A}_F/F}f\left(\begin{pmatrix}1&z\\ 0&1\end{pmatrix}x\right)dz=0$ 对所有 $x\in GL_2(\mathbb{A}_F)$ 成立；
\item[(3)] $f(xu)=f(x)$ 对所有 $x\in GL_2(\mathbb{A}_F)$、$u\in GL_2(\mathcal{O}_F)$ 成立。
\end{itemize}

用 $W$ 表示局部常值函数 $\phi:GL_2(\mathbb{A}_F)\to\overline{\mathbb{Q}}_\ell$ 构成的空间，使得
\begin{itemize}[leftmargin=2.2em,itemsep=2pt,topsep=2pt]
\item[(1$'$)] $\phi\left(\begin{pmatrix}a&z\\ 0&a\end{pmatrix}x\right)=\Psi(z)\phi(x)$ 对所有 $x\in GL_2(\mathbb{A}_F)$、$z\in\mathbb{A}_F$、$a\in F^\times$ 成立；
\item[(2$'$)] $f(xu)=f(x)$ 对所有 $x\in GL_2(\mathbb{A}_F)$、$u\in GL_2(\mathcal{O}_F)$ 成立。
\end{itemize}
% [原文疑误: 条件 (2$'$) 印作 $f(xu)=f(x)$，但此处讨论的是空间 $W$ 中的函数 $\phi$（条件 (1$'$) 亦用 $\phi$），应为 $\phi(xu)=\phi(x)$；系从上一段的 $V$ 的条件 (3) 照抄所致。译文照录原文。 来源：原书 p.124 英文原文与译本对照（A）。]

定义这两个空间之间的映射如下：
\begin{equation}\tag{11.1}
f\in V\mapsto\phi\in W;\qquad \phi(x)=\int_{\mathbb{A}_F/F}f\left(\begin{pmatrix}1&z\\ 0&1\end{pmatrix}x\right)\Psi(-z)\,dz,
\end{equation}
\begin{equation}\tag{11.2}
\phi\in W\mapsto f\in V;\qquad f(x)=\sum_{a\in F^\times}\phi\left(\begin{pmatrix}a&0\\ 0&1\end{pmatrix}x\right).
\end{equation}
根据 [JL]，这两个映射是互逆的同构。此外，它们与球 Hecke 代数在这两个空间上的（右）作用相交换。因此两个空间中的 Hecke 本征函数空间同构。相应的本征值由一族 $GL_2$ 共轭类 $\gamma_x$（$x\in C$）确定。

% ==== 原文 PDF 第 125 页（印刷页 125）====

已知对任意这样的族 $(\gamma_x)_{x\in C}$，$W$ 中的 Hecke 本征函数空间由所谓的 Whittaker 函数张成。为写出这个函数的显式公式 [W2]，我们用 $V_{m,k}$（$m\geq k$）表示以 $(m,k)$ 为最高权的 $GL_2$ 的不可约表示，即 $\mathrm{Sym}^{m-k}V\otimes(\det V)^{\otimes k}$，其中 $V$ 是定义二维表示。注意利用上述条件 (1$'$) 与 (2$'$)，$W$ 中的函数 $\phi$ 由它在形如
\begin{equation}\tag{11.3}
\left(\begin{pmatrix}t_x^{m_x}&0\\ 0&t_x^{k_x}\end{pmatrix}\right)_{x\in C}\in GL_2(\mathbb{A}_F)
\end{equation}
的元素上的值唯一确定（这里如前，$t_x$ 表示 $x$ 处的一致化子，即一个形式坐标）。给定共轭类的一族 $(\gamma_x)_{x\in X}$，$W$ 中相应的 Hecke 本征函数由公式
\begin{equation}\tag{11.4}
\phi\left(\left(\begin{pmatrix}t_x^{m_x}&0\\ 0&t_x^{k_x}\end{pmatrix}\right)\right)=\prod_{x\in C}\Tr(\gamma_x,V_{m_x+\delta_x,k_x})
\end{equation}
确定（相差一些无关紧要的非零因子）（若对某个 $x\in C$ 有 $m_x+\delta_x<k_x$，则按定义右端等于 0）。
% [Tutor Note: 推导补足。Whittaker 函数由它在 (11.3) 形元素上的值唯一确定，是因为条件 (1$'$) 允许把每个 $g\in GL_2(\mathbb{A}_F)$ 写成 $g=u\left(\begin{smallmatrix}a&z\\0&a\end{smallmatrix}\right)$（$u\in GL_2(\mathcal{O}_F)$，$a\in F^\times$，$z\in\mathbb{A}_F$）的乘积，而 $\Psi$ 与 $\phi$ 的变换性质把一般点上的值化为 (11.3) 形元素上的值。公式 (11.4) 中的指标平移 $\delta_x$ 来自微分 $\omega$ 的除子：$\Psi$ 的定义里的 $\operatorname{Res}_x(f_x\omega)$ 使中心特征标多出一个与 $\omega$ 的零点、极点有关的移位，故 $V$ 的最高权从 $(m_x,k_x)$ 变为 $(m_x+\delta_x,k_x)$。 来源：原书 p.125 上下文（A）。]

$GL_2$ 的、本征值对应于上述一族 $(\gamma_x)_{x\in X}$ 的尖点自守 Hecke 本征函数，按定义是满足上述条件 (2)、(3) 以及比 (1) 更强的条件的 Hecke 本征函数；这个更强的条件是：$f(\alpha x)=f(x)$ 对所有 $x\in GL_2(\mathbb{A}_F)$、$\alpha\in GL_2(F)$ 成立。换言之，它是双商
\begin{equation}\tag{11.5}
\Bun_{GL_2}(\mathbb{F}_q) = GL_2(F)\backslash GL_2(\mathbb{A}_F)/GL_2(\mathcal{O}_F)
\end{equation}
上的函数。上述结果表明：若它存在，则它必须由把变换 (11.2) 作用于具本征值 $(\gamma_x)_{x\in C}$ 的 Whittaker 函数得到。特别地，若存在，它在相差一个标量倍数意义下唯一。

根据 Drinfeld [Dr1, Dr2] 的结果，若 $\tilde\sigma: W_F\to GL_2$ 不可约且非分歧，则 (11.5) 上关于 $\tilde\sigma$ 的 Hecke 本征函数向量空间是一维的，且由尖点函数组成。因此这个空间的一个生成元由把算子 (11.2) 作用于具本征值 $(\tilde\sigma(\mathrm{Fr}_x))$ 的 Whittaker 函数给出。我们把这个自守函数记作 $f_{\tilde\sigma}$，相应的 Whittaker 函数记作 $W_{\tilde\sigma}$。

现在我们从 $GL_2$ 转向 $SL_2$。回忆我们想求 $\Bun^{\mathcal{O}(D)}_{SL_2}(\mathbb{F}_q)$ 上 Hecke 本征函数空间的维数，它是双商 (7.12)（这里如前，$D$ 是 $C$ 的有限子集，我们把它看作有效除子）。我们通过 $(g_x)\mapsto(\overline{g}_x)$ 把这个分量嵌入双商 (11.5)，其中
\[
\overline{g}_x = g_x\begin{pmatrix}t_x&0\\ 0&1\end{pmatrix},
\]
若 $x\in D$，而 $\overline{g}_x=g_x$，否则。我们将假定 $\sigma: W_F\to O_2\subset PGL_2$ 可以提升为同态 $\tilde\sigma: W_F\to GL_2$。容易看出，$\Bun^{\mathcal{O}(D)}_{SL_2}(\mathbb{F}_q)$ 上关于 $\sigma$ 的 Hecke 本征函数空间等于 (11.5) 上关于 $\tilde\sigma$ 的 Hecke 本征函数空间在此嵌入的像上的限制。

注意对上述形式的一般 $\sigma$，表示 $\tilde\sigma$ 将不可约。于是，要判定 $\Bun^{\mathcal{O}(D)}_{SL_2}(\mathbb{F}_q)$ 上关于 $\sigma$ 的 Hecke 本征函数空间是零维还是一维，

% ==== 原文 PDF 第 126 页（印刷页 126）====

我们需要判定上面构造的函数 $f_{\tilde\sigma}$ 在 $\Bun^{\mathcal{O}(D)}_{SL_2}(\mathbb{F}_q)$ 的像上的限制是否为零。这又由 Whittaker 函数 $W_{\tilde\sigma}$ 是否在所有形如 (11.3)、且除子
\[
\sum_{x\in C}(m_x+k_x)[x]
\]
线性等价于 $D$ 的元素上等于零所决定。

设 $\alpha: W_F\to\mathbb{Z}_2$ 是由 $\sigma$ 与同态 $O_2\to\mathbb{Z}_2$ 复合得到的二次特征标。它对应于一个二次扩张 $E/F$，我们在 §8.8 中称之为与 $\sigma$ 相伴的扩张。我们有
\[
\alpha(\mathrm{Fr}_x)=\begin{cases}1,&x\ \text{is split in }E,\\ -1,&x\ \text{is non-split in }E.\end{cases}
\]
对 $C$ 上任一除子 $M=\sum_{x\in C}M_x[x]$，令
\[
\langle M\rangle = \sum_{x\ \text{non-split in }E}M_x.
\]

回忆我们有用于定义特征标 $\Psi$ 的微分 $\omega$ 的除子 $\delta$，以及除子 $D=\sum_x D_x[x]$，其中 $D_x=0$ 或 $1$。

我们断言：Whittaker 函数 $W_{\tilde\sigma}$ 在所有形如 (11.3)、且
\[
\sum_{x\in C}(m_x+k_x)[x]
\]
线性等价于 $D$ 的元素上的值全为零，当且仅当 $\langle\delta+D\rangle$ 为奇数。事实上，这个集合包含元素 $(g_x)$，其中
\[
g_x=\begin{pmatrix}t_x&0\\ 0&1\end{pmatrix},\qquad x\in D,
\]
而 $g_x=1$，否则。按公式 (11.4)，$W_{\tilde\sigma}$ 在这个元素上的值等于
\[
\prod_{x\in C}\Tr(\tilde\sigma(\mathrm{Fr}_x),V_{\delta_x+D_x,0}).
\]
若 $\langle\delta+D\rangle$ 为奇数，则至少存在一个非分裂点 $y$，使得 $\delta_y+D_y$ 为奇数。但由于在此情形 $\tilde\sigma(\mathrm{Fr}_y)$ 共轭于矩阵 (7.5) 的一个标量倍数，我们发现
\[
\Tr(\tilde\sigma(\mathrm{Fr}_y),V_{\delta_y+D_y,0})=0.
\]
因此 $W_{\tilde\sigma}$ 在这个点处等于 0。我们需要检查的其他所有点都具有 (11.3) 的形式，且
\[
\sum_{x\in C}(m_x+k_x)[x]=D+(F),
\]
其中
\[
(F)=\sum_{x\in C}n_x[x]
\]

% ==== 原文 PDF 第 127 页（印刷页 127）====

是 $C$ 上有理函数 $F$ 的零点与极点的除子。但由交换类域论（例如见 [Se]），对 $C$ 上任一有理函数 $F$ 我们有
\begin{equation}\tag{11.6}
\prod_{x\in C}\alpha(\mathrm{Fr}_x)^{n_x}=\prod_{x\ \text{non-split in }E}(-1)^{n_x}=1.
\end{equation}
因此对满足上述条件的任一元素 (11.3)，同样至少存在一个非分裂点 $z\in C$，使得 $m_z+\delta_z-k_z$ 为奇数。公式 (11.4) 于是表明 $W_{\tilde\sigma}$ 在所有这些元素上的值为 0。因此我们发现，在此情形 $f_{\tilde\sigma}$ 在 $\Bun^{\mathcal{O}(D)}_{SL_2}(\mathbb{F}_q)$ 于 (11.5) 中的像上的限制等于 0。于是，$\Bun^{\mathcal{O}(D)}_{SL_2}(\mathbb{F}_q)$ 上关于 $\sigma$ 没有非零的 Hecke 本征函数，这正是我们要证明的。

另一方面，若 $\langle\delta+D\rangle$ 为偶数，则容易看出 $f_{\tilde\sigma}$ 在 $\Bun^{\mathcal{O}(D)}_{SL_2}(\mathbb{F}_q)$ 的像上的限制非零。因此 $\Bun^{\mathcal{O}(D)}_{SL_2}(\mathbb{F}_q)$ 上相应的 Hecke 本征函数空间是一维的。
% [Tutor Note: 推导补足。证明的核心是把两个对称性拼在一起：一方面，$\tilde\sigma$ 经二次特征标 $\alpha$ 约化到 $O_2$，故在非分裂点 $y$ 处 $\tilde\sigma(\mathrm{Fr}_y)$ 的射影类为 $\left(\begin{smallmatrix}0&1\\1&0\end{smallmatrix}\right)$，其在二维表示 $V_{\delta_y+D_y,0}$ 上的迹为零；另一方面，主除子 $(F)$ 上 $\alpha$ 的乘积由交换类域论等于 $1$（因为 $F$ 的典范类在 $\alpha$ 下平凡），这保证「$\langle\delta+D\rangle$ 为奇数」这一条件可以沿线性等价类传递，使所有相关元素上 $W_{\tilde\sigma}$ 都为零。这里 (11.6) 是整段论证的枢纽。 来源：原书 p.127 上下文（A）。]

为了完成定理 7.1 的证明，还须观察到
\[
\langle\delta\rangle = \sum_{x\ \text{non-split in }E}\delta_x
\]
总是偶数（见 [W1] 第 XIII 章第 12 节定理 13）。\footnote{我们感谢 B. Poonen 指出这一文献。}因此 $\langle\delta+D\rangle$ 的奇偶性与 $\langle D\rangle$ 相同，而后者与由 $D$ 的所有非分裂点构成的子集 $S\subset D$ 的元素个数相同。

% [以下为参考文献，属 references 分片]
